\subsection{Arbitrary relative kernels supported away from two and three}

\begin{statement}{E.20}{Proved}{Compatible extension through every relative quotient of order prime to six}
In this statement $N=N_B$ is the rank-$97$ kernel of
$B=C_2*C_3\to Q$, and hats denote full profinite completions.
Let $\mathscr P$ be the set of primes other than two and three.
For each open subgroup $L\triangleleft\widehat B$ contained
in $\widehat N$ with $[\widehat N:L]$ prime to six, and
integers $a,b\geq1$, put
\[
 K_{a,b,L}=
 \bigl(\overline{\gamma_3N\,N^{2^a3^b}}\bigr)\cap L.
 \tag{E.20.1}
\]
Every one of the $24$ admissible CM matrix pairs has a
compatible family of gentle shadows at the inverse images
of these subgroups in $\widehat B_3$, with multiplier one.
The family is compatible as $a,b$ increase and as $L$
decreases, and its projection to every quotient
$B/(\gamma_3N\,N^j)$ in E.15 is the family constructed
there. In particular it
includes simultaneous refinement by arbitrary relative
pro-$p$ quotients for every $p\geq5$, of every nilpotency
class, and by finite relative kernels whose orders involve
several primes in $\mathscr P$.

Here the three-strand action is constructed on one actual
profinite quotient, together with a conjugator in the image
of $\widehat F_2'$. No extra pentagon assertion is made.
The natural map to the finite quotient
\[
 B/\bigl(\gamma_4N\,(\gamma_2N)^3N^9\bigr)
 \tag{E.20.2}
\]
factors through none of the quotients in (E.20.1).
Its normal kernel over $Q$ is a finite three-group with
a nonzero third lower-central term, exhibited below.
Thus (E.20.1) supplies further compatible refinements while
leaving this specified relative three-primary comparison.
\end{statement}
\begin{proof}
Write $\mathcal N_{23}$ for the product of the free
class-two pro-two and pro-three completions of $N$ and
let $E_{23}$ be the corresponding quotient of $\widehat B$.
It fits into the exact sequence
\[
 1\longrightarrow\mathcal N_{23}\longrightarrow E_{23}
       \longrightarrow Q\longrightarrow1.
\]
Let $R_{\mathscr P}$ be the intersection of the kernels
of all continuous finite quotients of $\widehat N$
whose orders have prime factors only in $\mathscr P$.
This subgroup is characteristic in $\widehat N$, hence
normal in $\widehat B$. Put
\[
 E_{\mathscr P}=\widehat B/R_{\mathscr P},\qquad
 \mathcal N_{\mathscr P}=\widehat N/R_{\mathscr P}.
 \tag{E.20.3}
\]
Taking the $\widehat B$-core of a finite-index normal
subgroup of $\widehat N$ preserves this prime support:
its quotient embeds in a product of at most $|Q|$
conjugate quotients. Thus the subgroups $L$ in the
statement give the full topology of (E.20.3).

The simultaneous quotient is exactly the fibre product
\[
 H=E_{23}\mathbin{\times_Q}E_{\mathscr P}.
 \tag{E.20.4}
\]
To prove that the natural map $\widehat B\to H$ is
onto, first take finite quotients of the two normal
kernels. The image of $\widehat N$ in their product
surjects onto both factors. The two Goursat quotient
groups have a common order dividing a number supported
on $\{2,3\}$ and a number supported on $\mathscr P$.
That order is one. Hence the image of $\widehat N$
is the full product. Lifting the common $Q$ coordinate
then gives surjectivity on the entire finite fibre
product. The image of the compact group $\widehat B$
is closed; its surjectivity in every finite projection
proves (E.20.4).

The kernel $J$ of $\widehat B\to H$ satisfies
\[
                     J\subset\widehat N'\subset\widehat F_2'.
 \tag{E.20.5}
\]
Indeed the $\mathcal N_{23}$ coordinate detects every
two- and three-primary coordinate of
$\widehat N^{\rm ab}\simeq\widehat\Z^{97}$, and
$\mathcal N_{\mathscr P}$ detects every other primary
coordinate. All these coordinates vanish for an element
of $J$, proving the first inclusion. This also proves
the precise commutator comparison needed after gluing.
Let $H_F$ be the inverse image of $F_2/N$ under
$H\to Q$. There is a continuous homomorphism
\[
 H_F\longrightarrow\widehat F_2^{\rm ab},
 \qquad
 \ker(H_F\to\widehat F_2^{\rm ab})
       =\operatorname{im}(\widehat F_2'\to H_F).
 \tag{E.20.6}
\]
This is the map induced by abelianization, using
(E.20.5); the asserted equality follows by taking any
lift in $\widehat F_2$.

Fix one CM pair and use the ungauged two- and
three-primary automorphism $\alpha_{23}$ constructed
in E.15. Its action on $Q$ is the identity, and it
fixes the actual cusp element $t=sr$. Therefore
\[
 \alpha:H\longrightarrow H,
 \qquad (h_{23},h_{\mathscr P})\longmapsto
          (\alpha_{23}(h_{23}),h_{\mathscr P})
 \tag{E.20.7}
\]
is a continuous automorphism. It is invertible by the
same formula with $\alpha_{23}^{-1}$ and fixes $t$.
This construction uses the identity on the entire
relative $\mathscr P$ quotient, including all of its
commutator layers.

Take exactly the gauge $g\in\widehat\Z$ of E.15:
its two- and three-primary components are the gauges
there, and $g_p=0$ for every $p\in\mathscr P$.
With $x=t^{-2}$, define
\[
 \alpha'=\operatorname{Inn}(x^g)\circ\alpha,
 \qquad S'=\alpha'(s),\quad R'=\alpha'(r),\quad f=sS'.
 \tag{E.20.8}
\]
The gauge centralizes $t$, so
$S'^2=R'^3=1$ and $S'R'=t$. The word $f$ belongs
to $H_F$: the common $Q$ coordinate is precisely
$s x^{\bar g}s x^{-\bar g}$, as in E.15.11,
and its image in $B/F_2$ is the identity.

We check its full abelianization, rather than just
its reduction at two and three. Formula E.15.12 gives
\[
 \operatorname{ab}_{F_2}(f)
   =\operatorname{ab}_{F_2}(s\alpha(s))+(-g,g).
 \tag{E.20.9}
\]
The selected two- and three-primary gauges make
these components zero. For $p\in\mathscr P$, the
abelianization map to $\Z_p^2$ factors through
$E_{\mathscr P}$: its restriction to $\widehat N$
has finite $p$-group quotients, all retained in
(E.20.3). On that component $\alpha(s)=s$, so the
first term of (E.20.9) is zero, and $g_p=0$.
Consequently every primary component vanishes.
By (E.20.6), $f$ is the image of an actual element
of $\widehat F_2'$.

The torsion-to-hexagon comparison in E.9 applies to
these exact elements. In particular $\alpha'(x)=x$,
and, since $y=sxs$,
\[
             \alpha'(y)=S'xS'=f^{-1}yf.
 \tag{E.20.10}
\]
Both complete gentle hexagons hold with multiplier
one. The component automorphism $\alpha_{23}$
preserves every class-two finite level used in E.15;
the identity and its final inner conjugation preserve
every normal subgroup defining a quotient of
$E_{\mathscr P}$. Thus $\alpha'$ induces an
automorphism at every finite quotient in (E.20.1).
This proves generation, both equations, and all
transition compatibilities at once. Projection to
the class-two quotient of each primary kernel gives
the same ungauged actions and the same single gauge
as E.15. Hence the projected family is exactly the
one specified there and retains its CM pair.

We finish by proving the remaining finite comparison.
Set $L_3=\gamma_4N\,(\gamma_2N)^3N^9$, with the
closed interpretation in completions. Its quotient
in $N$ is finitely generated, has class at most
three and exponent dividing nine, and hence is a
finite three-group. Finiteness follows directly by
collecting words into the finitely many basic
commutators of weights at most three, each with
bounded exponent. The subgroup $L_3$ is characteristic
in $N$, so $B/L_3$ is an actual finite quotient.

The map from $N$ to $U_4(\F_3)$ given by
\[
 n_0\longmapsto I+E_{12},\quad
 n_1\longmapsto I+E_{23},\quad
 n_2\longmapsto I+E_{34},\quad
 n_i\longmapsto I\quad(i>2)
 \tag{E.20.11}
\]
kills $L_3$. The target has class three, its
commutator subgroup has exponent three, and
$(I+A)^3=I+A^3$ is central for a strictly upper
triangular $4\times4$ matrix $A$ over $\F_3$;
thus its exponent divides nine. Literal
multiplication gives
\[
 [[I+E_{12},I+E_{23}],I+E_{34}]=I+E_{14}\ne I.
 \tag{E.20.12}
\]
It follows that $[[n_0,n_1],n_2]$ has nonzero image
in $N/L_3$.

Suppose the natural quotient map $B\to B/L_3$
factored through one finite fibre product in
(E.20.1). The induced homomorphism to $Q$ agrees
with its given projection, because the maps agree
on the image of $B$. The normal kernel supported
on $\mathscr P$ therefore maps into $N/L_3$.
Its image has order prime to three and also
three-power order, so it is trivial. The map would
factor through the class-two $\{2,3\}$ component
and would kill $\gamma_3N$, contradicting
(E.20.12). This proves the claimed failure of
factorization by an exact quotient map.

The program \texttt{de\_coprime\_gluing\_root.py}
checks (E.20.11)--(E.20.12), enumerates all $729$
target elements, verifies their ninth powers and
central cubes, and verifies the indicated
commutator-subgroup exponent and central filtration.
The lifting and compatibility argument uses the
proved profinite morphisms (E.20.3)--(E.20.10).
\end{proof}
